\documentclass{article}
\usepackage{amsmath}
\usepackage{cancel}

\title{Degrees of Freedom}
\begin{document}
\maketitle

A \textbf{degree of freedom} is any independent physical parameter that
describes the state of a physical system. The set of all the possible states is
called the \textit{phase space} and the dimension of phase space is the number
of degrees of freedom. 

To put some numbers to this idea, consider space of dimension $D$ and a rigid
body moving through a $C\subseteq D$ subspace. Any rigid body can undergo two
distinct types of motion: \textit{translation} and \textit{rotation}. If there
are no constraints, then the body is free to move through all dimensions so that
the number of translational degrees of freedom is
\begin{equation}
  N_{\textrm{trans}} = D.
\end{equation}
This means that we would need at least $D$ generalized coordinates $q_i$ to
fully describe the system. The number of rotational degrees of freedom depends
on both the dimension of the total space and the dimension of the subspace
through which the object moves. In the case where $C=D$, then
\begin{equation}
  N_{\textrm{rot}} = \begin{pmatrix}D \\ 2\end{pmatrix} = \frac{1}{2}D(D-1)
\end{equation}
using the notation for binomial coefficients. The idea here is that rotation
occurs in planes and, therefore, the number of rotational degrees of freedom is
the number of independent planes we can slice through our $D$-dimensional space.

If instead we have that $C<D$, then the number of degrees of freedom is reduced
to
\begin{equation}
  N_{\textrm{rot}}= \begin{pmatrix}D\\2\end{pmatrix}-\begin{pmatrix}D-C\\2\end{pmatrix}
  = \frac{1}{2}D(D-1)-\frac{1}{2}(D-C)(D-C-1)
\end{equation}
because we have to subtract off the number of planes in the dimensions that the
object doesn't have access to. 


So far, we have
\begin{equation}
  N_{\textrm{total}} = N_{\textrm{trans}} + N_{\textrm{rot}}
\end{equation}


\subsection*{Constraints}
Often, it is easy to overdetermine the number of degrees of freedom. For
example: imagine a bead confined to move along rod in three dimensional space.
It would be easy to say that the bead has 3 translational degrees of freedom and
3 rotational degrees of freedom as we would expect for a free sphere moving
through space. In such a situation, we can still recover the correct number of
degrees of freedom by identifying all of the constraints on the system.
\textit{constraints} are just rules that prevent the bead from moving a certain
way. They often disguise complicated physics (normal forces, tension,
etc...) but can be written as a simple function of the form
\begin{equation}
  f(\vec{r}, t) = 0
\end{equation}
Constraints that can be written like this without depending on the coordinate
velocities are called \textit{holonomic constraints}. This distinction is
important because we will find that we can use Lagrangian mechanics to solve for
the forces that come from holonomic constraints. 


Returning to our main task, we now say that if there are are $K$ constraints,
the total degrees of freedom is given by
\begin{equation}
  N_{\textrm{total}} =  N_{\textrm{trans}} + N_{\textrm{rot}}-K
\end{equation}
Thus, if we overdetermine the system (perhaps we want to solve for the normal
force or something like that) then we can always recover the correct number of
degrees of freedom by subtracting off the number of constraint equations
relating the coordinates. 
\end{document}



